\scanpage{10}
le noyau de $A\longrightarrow B$ soit précisément $\unsure{\Gamma|H}$,
\deleted{\unsure{Pour|Prouvons} Montrer que $B$ est un schéma abélien}
\deleted{\unsure{Un}} Le passage au quotient en question
commute avec le changement de base, donc
$1^\circ$) $B\otimes_V K\simeq\deleted{\unsure{B}}A'/\Gamma'\simeq B'$ \qquad $2^\circ$) $B\otimes_V k$ est \deleted{\unsure{un}}
isom à \deleted{\unsure{un}} $A\otimes k/\Gamma'\otimes k$, donc est un
schéma abélien. Donc $B$ est un
schéma abélien. On a fini.

\uline{Corollaire} 1 \deleted{\unsure{Toute|Soit}} Si $A'$ \deleted{\unsure{est un}}\inserted{et $B'$ sont des}
schémas abéliens \inserted{isogènes} sur $K$, alors $A'$ ``se
réduit bien pour $V$'' sss $B'$ se réduit
bien pour $V$. S'il en est ainsi, les variétés
\deleted{\unsure{rédu}} abéliennes \unsure{réduites|induites} sont isogènes.

\deleted{\uline{Corollaire \unsure{2|3}}} Conséquence : On peut dire d'une
classe de variétés abéliennes mod isogénie
(p.ex. simple) si elle se réduit bien ou
non, et \deleted{\unsure{dire} \unsure{la}} parler de la
\deleted{variété \unsure{rédu}} classe réduite. Bien
entendu, une somme de classes
de \deleted{variétés}\inserted{schémas} abéliens sur $K$ qui se
réduisent bien, se réduit bien. En
particulier, si les composantes simples de
% archive sheet number: 9
